실험 플랫폼인 RP2040은 125 MHz 듀얼 Cortex-M0+와 264 KB SRAM으로 구성된
마이크로컨트롤러다\cite{rp2040}. 본 연구와 관련된 M0+의 특성은 세 가지다.

첫째, \textbf{분기가 싸다.} 2단 인오더 파이프라인이라 분기 예측기가 없으며, 분기는 방향과
무관하게 1--3사이클이다. 데스크톱 프로세서에서 15사이클 이상인 예측 실패 비용이 여기에는
존재하지 않는다.

둘째, \textbf{실행 시간이 결정적이다.} 대부분의 명령어가 1사이클, 로드가 2사이클로
고정되어\cite{m0plus}, 성능이 명령어 수로 직접 환산된다. 즉 ``몇 명령어로 짰는가''가 곧
성능이다.

셋째, \textbf{배럴 시프터에 포화 특성이 있다.} 레지스터로 시프트 양을 지정하는 명령어
(\texttt{LSLS Rd, Rm})는 Rm의 하위 8비트를 시프트 양으로 사용하고, 그 값이 32 이상이면
결과를 0으로 만든다\cite{armv6m}. 뒤에서 보듯 이 특성은 부호 판단을 하드웨어에 공짜로
떠넘길 수 있게 해 주지만, C 표준이 워드 폭 이상의 시프트를 미정의 동작(undefined
behavior)으로 규정하기 때문에 컴파일러는 이를 활용하는 코드를 생성할 수 없다.
